% Arch-1: the game-to-dataflow compilation pipeline, drawn as objects rather than
% labeled boxes. Single-column, drawn at native size and fitted to the column width.
% Top: an 11-node example tree (1 chance root, 2 P1 decisions, 2 P2 decisions,
% 6 terminals) over alternating depth bands, with one-word glyph labels at the
% first occurrence of each node kind. Middle: thick mapping lines funnel every
% node into its breadth-first cell of the flat node array, above the values
% template. Bottom: the run-time submission timelines (eager vs captured graph).
\begin{figure}[!htb]
\begin{minipage}[b]{0.50\linewidth}\centering
\resizebox{\linewidth}{!}{%
\begin{tikzpicture}[
    font=\sffamily\small,
    dec1/.style={circle, draw=teal!60!black, fill=teal!45, minimum size=3.4mm, inner sep=0pt},
    dec2/.style={circle, draw=violet!60!black, fill=violet!40, minimum size=3.4mm, inner sep=0pt},
    chn/.style={diamond, draw=black!65, fill=black!22, minimum size=4.2mm, inner sep=0pt},
    ter/.style={rectangle, draw=black!70, fill=black!55, minimum size=2.5mm, inner sep=0pt},
    edge/.style={black!55, line width=0.9pt},
    flow/.style={-{Stealth[length=2.2mm]}, line width=1.0pt, black!60},
    fold/.style={-{Stealth[length=1.8mm]}, densely dashed, line width=0.8pt, black!55},
    mapline/.style={line width=1.0pt, opacity=0.80, shorten <=1.4pt},
    tick/.style={orange!80!black, line width=1.0pt},
    gtick/.style={orange!95!black, line width=1.6pt},
    note/.style={font=\sffamily\small, text=black!55},
    lanelabel/.style={font=\sffamily\small\bfseries, text=black!60}
]
% ================= build lane: tree over its flattened arrays =================
% alternating depth bands behind the tree
\foreach \y in {0, -1.7}
  \fill[blue!6] (0.10,\y+0.33) rectangle (7.95,\y-0.33);
% --- example game tree; left margin below stays clear for the compile label ---
\node[chn]  (c0)  at (4.60, 0.00) {};
\node[dec1] (p1a) at (2.60,-0.85) {};
\node[dec1] (p1b) at (6.60,-0.85) {};
\node[dec2] (q2a) at (1.70,-1.70) {};
\node[ter]  (t2a) at (3.50,-1.70) {};
\node[dec2] (q2b) at (5.70,-1.70) {};
\node[ter]  (t2b) at (7.60,-1.70) {};
\node[ter]  (t3a) at (1.30,-2.55) {};
\node[ter]  (t3b) at (2.30,-2.55) {};
\node[ter]  (t3c) at (5.20,-2.55) {};
\node[ter]  (t3d) at (6.20,-2.55) {};
\draw[edge, densely dashed] (c0) -- (p1a);
\draw[edge, densely dashed] (c0) -- (p1b);
\draw[edge] (p1a) -- (q2a); \draw[edge] (p1a) -- (t2a);
\draw[edge] (p1b) -- (q2b); \draw[edge] (p1b) -- (t2b);
\draw[edge] (q2a) -- (t3a); \draw[edge] (q2a) -- (t3b);
\draw[edge] (q2b) -- (t3c); \draw[edge] (q2b) -- (t3d);
% a compact standalone glyph legend keeps labels off the tree
\node[chn, minimum size=3.2mm] (legchance) at (0.55,0.62) {};
\node[note, anchor=west] at (0.98,0.62) {chance};
\node[dec1, minimum size=3.0mm] (legp1) at (2.30,0.62) {};
\node[note, anchor=west] at (2.73,0.62) {P1};
\node[dec2, minimum size=3.0mm] (legp2) at (3.48,0.62) {};
\node[note, anchor=west] at (3.91,0.62) {P2};
\node[ter, minimum size=2.4mm] (legterm) at (4.70,0.62) {};
\node[note, anchor=west] at (5.13,0.62) {terminal node};

% ================= flat arrays (BFS order, depth-segmented) =================
\def\cw{0.62}
\def\sx{0.72}
\def\slotb{-4.17}   % node strip bottom (top = -4.05)
\def\tmplb{-5.87}   % template strip bottom (top = -5.75)
% mapping fan: every tree node funnels into its BFS cell (compile = walk once).
% Each line first drops straight out of the tree, then converges inside the
% dedicated fan band so crossings never overlap the tree itself.
\foreach \n/\i/\c in {c0/0/black!50, p1a/1/teal!75, p1b/2/teal!75,
                      q2a/3/violet!75, t2a/4/black!50, q2b/5/violet!75,
                      t2b/6/black!50, t3a/7/black!50, t3b/8/black!50,
                      t3c/9/black!50, t3d/10/black!50}
  \draw[mapline, \c] (\n.south) -- (\n.south |- 0,-2.85)
    .. controls (\n.south |- 0,-3.25) and ({\sx+(\i+0.5)*\cw},-3.10) ..
    ({\sx+(\i+0.5)*\cw},\slotb+\cw);
\node[note, align=left, anchor=west] at (0.00,-3.02) {compile\\walk once};
% node strip: cells keep the color of their tree node, in BFS order
\node[note, anchor=east] at (\sx-0.10,\slotb+0.5*\cw) {nodes};
\foreach [count=\i from 0] \c in
  {black!22, teal!45, teal!45, violet!40, black!55, violet!40, black!55,
   black!55, black!55, black!55, black!55}{
  \draw[black!60, fill=\c] ({\sx+\i*\cw},\slotb) rectangle ++(\cw,\cw);
}
% chance slot carries no run-time work: overlay dashed hatch look
\draw[black!60, densely dashed, fill=white, fill opacity=0.55]
  (\sx,\slotb) rectangle ++(\cw,\cw);
% depth segment separators + labels
\foreach \s in {0, 1, 3, 7, 11}
  \draw[black!85, line width=1.2pt] ({\sx+\s*\cw},\slotb-0.05) -- ++(0,\cw+0.10);
\foreach \d/\a/\b in {0/0/1, 1/1/3, 2/3/7, 3/7/11}
  \node[note] at ({\sx+(\a+\b)*\cw/2},\slotb-0.26) {$d{=}\d$};
% per-depth index blocks: one gather/scatter block arrow per segment
\foreach \a/\b in {1/3, 3/7, 7/11}
  \draw[flow, line width=0.9pt, -{Stealth[length=1.8mm]}]
    ({\sx+\a*\cw+0.06},\slotb-0.56) -- ({\sx+\b*\cw-0.06},\slotb-0.56);
\node[note] at ({\sx+9*\cw},\slotb-0.86) {$D$ blocks};
% values template strip: only terminal cells are pre-filled
\node[note, anchor=east] at (\sx-0.10,\tmplb+0.5*\cw) {$\cfv_{\mathrm{tmpl}}$};
\foreach [count=\i from 0] \c in
  {white, white, white, white, black!55, white, black!55,
   black!55, black!55, black!55, black!55}{
  \draw[black!60, fill=\c] ({\sx+\i*\cw},\tmplb) rectangle ++(\cw,\cw);
}
\foreach \s in {0, 1, 3, 7, 11}
  \draw[black!85, line width=1.2pt] ({\sx+\s*\cw},\tmplb-0.05) -- ++(0,\cw+0.10);
% chance folding: the run-time-free chance slot melts into the template
\draw[fold] ({\sx+0.04},\slotb)
  .. controls (0.28,-4.60) and (0.28,-4.95) .. ({\sx+0.08},\tmplb+\cw)
  node[note, pos=0.5, left=0.6mm] {fold};

% ================= run lane: submission timelines =================
\begin{scope}[yshift=-0.10cm]
\def\wzero{1.66}
\def\itw{1.62}
\def\itg{0.14}
% three iteration windows, shaded alternately, shared by both rows
\foreach \i/\f in {0/black!5, 1/black!0, 2/black!5}
  \fill[\f] ({\wzero+\i*(\itw+\itg)},-7.30) rectangle ++(\itw,-1.15);
\node[note] at ({\wzero+0.5*\itw},-7.46) {one iteration};
% eager row: one launch tick per tensor op, dense within every iteration
\node[note, anchor=east] at (0.72,-7.78) {eager};
\draw[black!35] (\wzero,-7.78) -- ({\wzero+3*\itw+2*\itg},-7.78);
\foreach \i in {0,1,2}{
  \foreach \k in {0,...,8}{
    \draw[tick] ({\wzero+0.18+\i*(\itw+\itg)+\k*0.17},-7.65) -- ++(0,-0.26);
  }
}
% graph row: record once, then a single submission per iteration
\node[note, anchor=east] at (0.72,-8.25) {graph};
\draw[black!35] (\wzero,-8.25) -- ({\wzero+3*\itw+2*\itg},-8.25);
\node[draw=orange!80!black, fill=orange!25, rounded corners=1.5pt,
      minimum height=3.4mm, inner sep=1.6pt, anchor=east]
  at ({\wzero-0.08},-8.25) {record};
\foreach \i in {0,1,2}{
  \draw[gtick] ({\wzero+0.18+\i*(\itw+\itg)},-8.06) -- ++(0,-0.38);
}
\draw[flow, line width=0.9pt]
  ({\wzero+3*\itw+2*\itg},-8.25) -- ++(0.26,0);
\node[note, align=left, anchor=west] at (7.10,-8.25) {replay\\$\times T$};
% time axis hint
\draw[-{Stealth[length=1.6mm]}, black!40] (5.50,-8.65) -- (6.70,-8.65)
  node[note, pos=0.5, above=0.2mm] {time};
\end{scope}

% ================= lane frames =================
\begin{scope}[on background layer]
\node[draw=blue!45!black!30, rounded corners=3pt, inner sep=1.2mm, fill=blue!3,
      fit={(-0.22,0.88) (8.14,-6.35)}] (buildlane) {};
\node[draw=orange!70!black!30, rounded corners=3pt, inner sep=1.2mm, fill=orange!4,
      fit={(-0.22,-7.15) (8.14,-9.00)}] (runlane) {};
\end{scope}
\node[lanelabel, anchor=south west] at (buildlane.north west)
  {Build: once per game};
\node[lanelabel, anchor=south west] at (runlane.north west)
  {Run: every iteration};
\end{tikzpicture}}
\caption{The compilation pipeline on an 11-node example game. A build-time walk
flattens the tree into the depth-segmented node array and values template
(colored lines trace each node to its cell); at run time, eager execution
submits $c_1 + c_2 D$ framework operations per iteration (each one or more kernel
launches), graph replay one graph launch.}
\Description{Two-lane diagram. The build-time lane shows a small game tree with
a diamond chance root, colored decision nodes, and square terminals over
alternating depth bands; thick colored lines funnel every node into its
breadth-first cell of a flat array strip divided into depth segments, above a
values-template strip whose terminal cells are pre-filled; a short dashed arrow
folds the chance slot into that template. The run-time lane shows two aligned
timelines over three shaded iteration windows: the eager row has a dense comb
of launch ticks in every window, while the graph row has one record block
followed by a single tick per window and an arrow labeled replay times T.}
\label{fig:archpipeline}
\end{minipage}\hfill%
